\par\SourceLeaf{134}
Gerade hiedurch erhält nun aber die Unklarheit\footnote{Wer an dieser Bezeichnung einer seit einem halben Jahrhundert in Kraft gewesenen Anordnung Anstoss nimmt, wolle sein Urtheil suspendiren, bis er die erste Abtheilung ganz gelesen hat.} des Regulativs bei derjenigen Stipulation, wodurch die Grösse und das Fortschreiten des Pensionssatzes normirt werden soll, einen sehr bedenklichen Charakter. Die eignen Worte des Rescripts des Universitätscuratoriums vom 20. November 1794, durch welche diese Normirung sanctionirt ist, sind folgende:
\begin{quote}
— — — — Zweitens genehmigen wir, dass so oft sich der Fundus um 5000 Thl. vermehrt haben wird, und die andern Revenüen des Witwenfonds keine Verminderung gelitten haben, auch die Anzahl der Pensionen nicht über 15 gestiegen ist, eine jede Pension mit 10 Thl. erhöhet werden solle.
\end{quote}
\noindent\hspace*{\parindent}Mangelhaft ist diese Bestimmung darin, dass sich nicht auf eine ganz unzweideutige Art erkennen lässt, \textit{was} denn eigentlich an einer Pensionserhöhung von den vorangehenden Bedingungen abhängig sein soll, ob das Bestehen, oder ob bloss der Anfang; mit andern Worten (indem ich bloss die letzte Bedingung in Betracht ziehe)
\begin{quote}
ob eine Erhöhung, die dem Statut gemäss zu einer Zeit eingetreten ist, wo die Zahl der Witwen höchstens 15 betragen hatte, wieder aufhören oder wenigstens zweckmässig modificirt werden soll, sobald später die Anzahl der Pensionen jene Normalzahl überschreitet,
\end{quote}
oder aber
\begin{quote}
ob eine einmal eingetretene Erhöhung unabhängig von der späterhin erfolgenden Ueberschreitung der Normalzahl dennoch unabänderlich fortdauern solle.
\end{quote}
\noindent\hspace*{\parindent}Dass die obige Formulirung der Vorschrift, wenn man ohne alle Rücksicht auf die bei der Erwägung des Inhalts sich ergebenden Folgerungen, bloss den Wortlaut in Betracht zieht, natürlicher auf die zweite Auslegung hinführt als auf die erste, will ich um so weniger bestreiten, da bei meiner gegenwärtigen hauptsächlich auf den innern Gehalt der Anordnung selbst gerichteten Untersuchung der sprachliche Standpunkt nur ein ganz untergeordneter ist. Ich kann jedoch nicht umhin, zur Vergleichung auch die Einkleidung hieher zu setzen, in welcher \textsc{Brandes}, in seinem bekannten Werke über die Universität Göttingen, die Verfügung anführt; da \textsc{Brandes} 1794 als Referent im Ministerium für die Universitätssachen fungirte, so ist seine Auffassung jedenfalls zur Sache gehörig, wenn sie auch vom juristischen (meiner Untersuchung gleichfalls fremden) Standpunkt aus, der einmal im officiellen Rescript gebrauchten Wortfassung nicht derogiren kann. Es heisst a. a. O. S. 254:
\begin{quote}
Ferner wurde zu der Zeit beliebt, dass jedesmal, wenn der Capital-Fonds der Kasse mit 5000 Thl. angewachsen sei, eine jede Pension mit 10 Thl. vermehrt werden solle, \textit{so lange} die Zahl der Pensionirten nicht über 15 hinausgeht, was noch nie der Fall war.
\end{quote}
\noindent\hspace*{\parindent}Hiernach scheint \textsc{Brandes} den Vorbehalt wegen der 15 Pensionen eher in dem Sinn der \textit{ersten} Interpretation verstanden, aber, wegen des zuletzt von ihm angeführten Umstands für praktisch ganz unerheblich gehalten zu haben, wodurch sich denn vielleicht auch erklären lässt, dass in der Wortfassung des Rescripts nicht für die vollkommenste Schärfe Sorge getragen ist.

\textsc{Brandes} würde jedoch wahrscheinlich ganz anders geurtheilt haben, wenn er vorher das Factische genau geprüft hätte. \textsc{Brandes} kann bei der Angabe ‘was noch nie der Fall war’ nicht den siebenjährigen Zeitraum von Einführung der Bestimmung bis zur Abfassung seines Buchs verstanden haben, da der Erfahrung aus einem so kurzen Zeitraume gar kein Gewicht beigelegt werden könnte, sondern muss die ganze seit Stiftung der Kasse verflossene Zeit gemeint haben. Dann ist aber seine Behauptung factisch unrichtig. Die Zahl der aus der Kasse bezahlten Pensionen war wirklich schon einmal über 15 gestiegen, nemlich während der ersten sechs Monate des Jahrs 1778, wo 16 Pensionen bestanden; ja derselbe Fall würde
